\chapter{Theory and Frameworks}

This chapter provides the theoretical foundation for understanding the methods in Chapter 3. It covers scRNA-seq preprocessing concepts, the COTAN framework for sparse data, and Differential Transcript Usage fundamentals.

\section{Preprocessing Pipeline}

This section follows a single scRNA-seq experiment from its data-generating mechanism to the count matrix consumed by every later analysis, and treats that path as a four-phase pipeline: ingestion and quality control, mapping, probabilistic resolution of ambiguous reads, and deduplication into a matrix. Each phase is constrained by the one before it, so the hardware chosen in the first step ultimately dictates the shape, sparsity, and content of the final matrix.

\subsection{The Data-Generating Mechanism: Throughput vs.\ Resolution}
scRNA-seq protocols fall into two categories with fundamentally different characteristics:

\begin{itemize}
    \item \textbf{Full-length protocols (SMART-seq2\defn{SMART-seq2}{A plate-based full-length scRNA-seq protocol, one cell per well, capturing whole transcripts.}):} Capture entire transcripts from 5' to 3', enabling comprehensive isoform resolution. These use plate-based chemistry, producing higher per-cell read depth but limiting throughput to hundreds or thousands of cells.
    
    \item \textbf{$3'$-biased droplet methods (10x Genomics\defn{10x Genomics Chromium}{The dominant commercial droplet platform, barcoding each cell in a microfluidic droplet.}):} The widely used Chromium platform sequences only the extreme transcript termini using microfluidic barcoding. This achieves massive throughput (tens to hundreds of thousands of cells) but loses internal splice junction information critical for isoform distinction.
\end{itemize}

The throughput gap follows directly from how cells are compartmentalised. In plate-based protocols, each cell occupies one well of a 96- or 384-well plate, so cells must be dispensed individually by fluorescence-activated cell sorting\defn{Fluorescence-activated cell sorting (FACS)}{A method that sorts single cells into wells one at a time using fluorescent labels.} or limiting dilution, and every well constitutes a separate reaction requiring its own reagent volumes and handling steps. Throughput therefore scales with the number of available wells and the manual or robotic operations per plate, which caps a typical experiment at hundreds to a few thousand cells. Droplet microfluidics\defn{Droplet microfluidics}{A technique that forms thousands of picolitre droplets per second to compartmentalise individual cells.} instead encapsulates each cell together with a barcoded bead inside a picolitre droplet, forming thousands of droplets per second in a single microfluidic channel. The same instrument can thus process tens of thousands of cells in one run, at the cost of the reduced per-cell read depth imposed by sharing the sequencer's capacity across a far larger population.

This trade-off is a hardware constraint, not a biological one, and it propagates through the entire pipeline. Droplet methods deliver matrices with a large number of cells $N$ but shallow per-cell depth, which produces the zero-inflated structure that dominates the remainder of this chapter. Plate-based methods deliver fewer cells at higher depth, recovering internal splice junctions at the cost of scale.

\subsection{Phase 1: Data Ingestion and Quality Control}
Before quantification, raw sequencing reads undergo quality control. Adapter sequences are trimmed and low-quality bases filtered using tools such as cutadapt or fastp. For droplet-based protocols (e.g., 10x Genomics), cell barcode\defn{Cell barcode}{A short DNA tag in every read that identifies the cell of origin, allowing reads to be assigned to their cell.} filtering is applied to exclude empty droplets and multiplets.

Droplet protocols additionally carry two barcodes on every read: a cell barcode identifying the cell of origin, and a Unique Molecular Identifier (UMI)\defn{Unique Molecular Identifier (UMI)}{A random barcode attached to each RNA molecule before PCR amplification. It lets the pipeline tell true biological copies apart from PCR artifacts, preventing overcounting of amplified transcripts.}---a random sequence ligated to each captured mRNA before PCR amplification. FASTQ\defn{FASTQ}{The file format holding raw sequencing reads plus a per-base quality score for each nucleotide.} preparation extracts and trims these tags.


\subsection{Phase 2: Mapping Algorithms (Alignment vs.\ Pseudo-Alignment)}
Quantifying transcript expression requires assigning each short sequencing read to the transcript it originated from. The total computational workload is governed by the number of reads,

$$R = N \cdot M,$$

where $N$ is the number of cells and $M$ is the average number of reads per cell. In a large-scale droplet experiment, $N \approx 5 \times 10^4$ cells and $M \approx 2 \times 10^6$ reads per cell, so $R \approx 10^{11}$ reads.

The remaining dimensions follow from the mammalian references and the $10$x Genomics chemistry analysed in this thesis (Chapter~3). The transcriptome holds $T \approx 200{,}000$ isoforms, the order of magnitude of a standard human or mouse GENCODE annotation (roughly 250{,}000 transcripts in human and 150{,}000 in mouse). Because unspliced pre-mRNA is filtered out during preprocessing, the pseudo-alignment index contains only mature, spliced transcripts, whose average length is $L_{avg} \approx 1{,}500$\,bp. Read length is a hardware constraint rather than a biological one: in a $3'$ assay, Read~1 carries the cell barcode and UMI while Read~2 captures the transcript sequence, sequenced at $L \approx 90$--$150$\,bp depending on the Illumina kit chemistry.

Traditional alignment algorithms\defn{Read alignment}{Base-by-base mapping of a read to the reference, giving exact coordinates at high computational cost.} (e.g., BWA, Bowtie2) map reads base-by-base against the reference using a Burrows-Wheeler Transform\defn{Burrows-Wheeler Transform (BWT)}{A reversible text transform that allows fast substring search in a compressed reference index.} index. Constructing the index costs $O(T \cdot L_{avg})$, and each of the $R$ reads is then located by a search whose cost grows logarithmically with the size of the reference. The query phase therefore costs $O(R \cdot L \log T)$, giving a total execution time of

$$O(T \cdot L_{avg}) + O(R \cdot L \log T).$$

Each read is effectively searched against the entire transcriptome, and the aligner must additionally tolerate sequencing errors, insertions and deletions. This yields exact base-level coordinates, but the cost of the search grows with the size of the reference, which becomes prohibitive across the billions of reads in a single-cell experiment.

Modern droplet-based transcriptomics relies on \textbf{pseudo-alignment}, which rests on a narrower observation: quantification does not require a read's exact coordinates, only the identity of the transcript it came from. Pseudo-alignment algorithms (such as Salmon/alevin-fry or kallisto) therefore replace the search with table lookups. Each read is split into overlapping $k$-mers\defn{k-mer}{A contiguous substring of length $k$; reads are matched to transcripts by the $k$-mers they share.}---short substrings of length $k$---and every $k$-mer is looked up in a table recording which transcripts contain it. Intersecting these sets leaves the transcripts compatible with the read, an equivalence class\defn{Equivalence class}{The set of transcripts that could have produced a read, since all contain the read's $k$-mers.} of candidate transcripts\cite{he2022alevin}. The algorithm proceeds in two phases:

\begin{itemize}
    \item \textbf{Index construction:} a $k$-mer lookup table (or colored de Bruijn graph\defn{Colored de Bruijn graph}{A graph of overlapping $k$-mers in which each $k$-mer is coloured by the transcripts containing it, used as the pseudo-alignment index.}, in which each $k$-mer is coloured by the transcripts that contain it) is built once from the reference transcripts, at a cost of
    $$O(T \cdot L_{avg}) \;\text{(time)}, \qquad O(k \cdot |\Sigma|^k) \;\text{(space)},$$
    where $|\Sigma|=4$ for DNA. This cost is identical for both alignment and pseudo-alignment, and is paid once before any read is processed.
    
    \item \textbf{Query phase:} for each read, all $L-k+1$ overlapping $k$-mers are extracted and their transcript sets intersected. Each step is a hash-table lookup, so the per-read cost is independent of transcriptome size:
    $$O(L).$$ With $k=31$ and $L \approx 90$--$150$\,bp, roughly 60--120 lookups are performed per read.
\end{itemize}

The two methods therefore share the identical index-construction term $O(T \cdot L_{avg})$ and differ only in the query term. Pseudo-alignment replaces the logarithmic search $O(R \cdot L \log T)$ with a constant-time lookup $O(R \cdot L)$, yielding a total execution time of

$$O(T \cdot L_{avg}) + O(R \cdot L).$$

Decoupling the per-read cost from the size of the reference is what produces the speedup. The asymptotic gap between the two query terms is a factor of $\log_2 T \approx 17.6$ for $T = 200{,}000$, and the lower constant factor of a hash lookup widens the observed speedup to the 50--100× range typically reported. For the droplet experiment above ($N = 5 \times 10^4$, $M = 2 \times 10^6$, $R \approx 10^{11}$), pseudo-alignment reduces processing time from days to hours on standard hardware.


\subsection{Phase 3: Probabilistic Resolution via Expectation-Maximization}
Pseudo-alignment introduces a computational challenge: \textbf{multi-mapping reads}\defn{Multi-mapping read}{A read compatible with more than one transcript because those transcripts share its sequence.}---reads whose $k$-mers match more than one transcript because those transcripts share exonic sequence. A read falling entirely within a shared exon carries no information about which isoform it came from, so pseudo-alignment returns a set of candidates rather than a single transcript. This ambiguity is most severe for $3'$-biased protocols, where short reads capture only the terminal region and often match every isoform that shares that exon \cite{galfre2021cotan}.

Discarding ambiguous reads is not a viable solution: they form the majority in $3'$-biased data, so dropping them would shrink the count matrix and bias isoform proportions toward whichever isoforms happen to carry unique sequence. Modern pseudo-aligners instead resolve the ambiguity with the Expectation-Maximization (EM)\defn{Expectation-Maximization (EM)}{An iterative algorithm that estimates transcript abundances and resolves ambiguous read assignments jointly until they converge.} algorithm, which estimates isoform abundances and read assignments jointly by alternating two steps:

\begin{itemize}
    \item \textbf{E-step}\defn{E-step / M-step}{The two alternating EM phases: distribute reads by the current abundance estimate (E), then re-estimate abundances (M).}: given the current abundance estimate, each multi-mapping read is distributed across its candidate transcripts in proportion to those abundances. A read compatible with two isoforms at 10 and 1 copies is assigned roughly $10/11$ and $1/11$ of a count.
    \item \textbf{M-step:} transcript abundances are re-estimated from these fractional assignments.
\end{itemize}

The two steps alternate until the estimates converge. This probabilistic split preserves relative isoform signals instead of discarding them\cite{he2022alevin}, which matters for DTU: the signal that distinguishes similar isoforms is exactly the signal an ambiguous read carries. EM therefore yields fractional counts, which are summed into the count matrix and rounded where integer counts are required.

\subsection{Phase 4: Deduplication and Matrix Generation}
Once pseudo-alignment has assigned each read to a transcript, reads sharing a cell barcode and a UMI at the same feature are collapsed into a single count. The collapse is applied per feature: a UMI identifies a molecule only within one cell and one gene, since independent molecules can coincidentally share the same UMI sequence. Feature assignment must therefore precede deduplication, as collapsing on the raw (cell, UMI) pair alone would merge distinct molecules. This defines the unit of quantification as the captured molecule rather than the sequencing read, so that PCR-amplified copies of one molecule are not mistaken for independent observations and cells sequenced at different depths remain comparable.

The preprocessing pipeline produces an expression count matrix\defn{Count matrix}{A table with rows indexing genes or transcripts, columns indexing cells, and each entry recording the observed read count for that feature in that cell.} $X \in \mathbb{N}^{G \times N}$, where rows index genomic features (genes or transcripts), columns represent distinct single cells, and entry $x_{fc}$ records the observed read count for feature $f$ in cell $c$. Throughout this thesis, $N$ denotes the number of cells and $G$ the number of features. Transcript-level resolution preserves isoform-specific signals that gene-level aggregation would obscure.


For downstream DTU analysis with COTAN, the matrix must be transcript-level rather than gene-level. Additionally, EM algorithms produce non-integer fractional counts which may require rounding before binary discretization in \gls{contingencyTable} analysis.

\section{COTAN Framework}

The preceding section produced a transcript-level matrix $X \in \mathbb{N}^{G \times N}$ by way of the four pipeline phases. That matrix is the natural input for the statistical framework presented here, and it is also the source of the difficulty: the shallow per-cell depth imposed by the droplet hardware leaves most of its entries at zero. The framework was introduced by Galfr\`e et al.\ \cite{galfre2021cotan} and has since been extended to recover cell clusters directly from co-expression structure \cite{galfre2026gene}. The following subsections characterise that sparsity and present the contingency-table machinery COTAN uses to analyse it directly.

\subsection{The Sparsity Problem}
Single-cell count matrices exhibit extreme sparsity: 70--90\% of entries are zeros\cite{galfre2021cotan}. These represent two distinct phenomena:
\begin{itemize}
    \item \textbf{Biological zeros}\defn{Biological zero}{A zero count reflecting genuine absence of a transcript in a cell, as opposed to a technical dropout.}: True absences where a gene/transcript is not transcribed in a specific cell state.
    \item \textbf{Technical zeros (dropouts):} False absences from failed capture, reverse-transcription, or sequencing due to low chemical efficiency (~10--20\%); they are indistinguishable from true biological zeros in raw counts.
\end{itemize}

Biological and technical zeros are numerically indistinguishable in raw count matrices.


This poses a fundamental statistical challenge for methods designed for dense bulk RNA-seq data:

\begin{enumerate}
    \item \textbf{Log-normalization\defn{Log-normalization (pseudo-count)}{Adding a small constant then taking the log to stabilise variance, which distorts sparse low counts.} artifacts:} Adding pseudo-counts then logging distorts low-abundance measurements, compressing zeros into artificial scaling patterns \cite{galfre2021cotan}.
    
    \item \textbf{Correlation distortion:} Pearson/Spearman correlations\defn{Pearson / Spearman correlation}{Standard linear and rank-based association measures that, on sparse matrices, track shared zeros rather than true co-expression.} measure dropout coincidence rather than true regulatory co-expression in sparse matrices\cite{galfre2021cotan}.
    
    \item \textbf{Imputation risks:} Heuristic smoothing replaces zeros with estimated values based on neighboring cells, frequently introducing artificial correlations and spurious downstream signals \cite{galfre2021cotan}.
\end{enumerate}

While some single-cell specific methods have been developed to address these issues, many still rely on imputation or transformation that can obscure true biological signal.


COTAN takes a different approach—treating zeros explicitly through binary discretization without requiring imputation or variance-stabilizing transformations.

\subsection{Contingency Table Solution}
COTAN bypasses continuous count amplitudes, evaluating pairwise dependencies using asymptotic independence testing\defn{Asymptotic independence test}{Using the large-sample $\chi^2$ distribution to test whether two features are statistically independent.} on binary presence/absence states\cite{galfre2021cotan}. For a feature pair $(f, f')$ across $N$ cells, each cell $c$ is reduced to the pair of indicators

$$a_c = \mathbf{1}[x_{fc} > 0], \qquad b_c = \mathbf{1}[x_{f'c} > 0],$$

and the states are indexed by $i, j \in \{0,1\}$ throughout this section. Let $O_{ij}$ denote the observed number of cells with $a_c = i$ and $b_c = j$, that is, the entries of the $2 \times 2$ contingency table. The marginal totals are $\mathrm{Row}_i = \sum_j O_{ij}$ and $\mathrm{Col}_j = \sum_i O_{ij}$, with $\sum_i \mathrm{Row}_i = \sum_j \mathrm{Col}_j = N$. Under the null hypothesis of independence, the expected count in table cell $(i,j)$ is

$$E_{ij} = \frac{\mathrm{Row}_i \cdot \mathrm{Col}_j}{N}.$$

Deviation from independence is evaluated via Pearson's chi-squared statistic\defn{Pearson's chi-squared statistic ($\chi^2$)}{A test statistic comparing observed counts with those expected under independence in a contingency table.}\cite{galfre2021cotan}:

$$\chi^2 = \sum_{i \in \{0,1\}} \sum_{j \in \{0,1\}} \frac{(O_{ij} - E_{ij})^2}{E_{ij}}.$$

With thousands of cells ($N \gg 30$), the asymptotic approximation remains robust\cite{galfre2021cotan}, enabling efficient testing across millions of feature pairs.

\subsection{Output Metrics}
COTAN generates three primary metrics \cite{galfre2021cotan}:
\begin{enumerate}
    \item \textbf{Co-expression matrix (coex)\defn{Co-expression matrix (coex)}{COTAN's normalized pairwise associations computed from contingency statistics. Negative values indicate mutual exclusivity (two features rarely co-occur).}:} Normalized pairwise associations from contingency statistics. Negative values indicate mutual exclusivity.
    
    \item \textbf{Differential Expression Analysis (DEA)\defn{DEA enrichment coefficient ($D_{f,S}$)}{A score in $(-1,1)$ measuring how strongly feature $f$ is enriched (positive) or depleted (negative) in cell subset $S$.}:} Enrichment coefficients $D_{f,S} \in (-1, 1)$ measuring feature--cluster association strength.
    
    \item \textbf{\gls{gdi}:} Cluster specificity measure for marker identification and clustering quality validation.
\end{enumerate}

Negative co-expression between isoforms of the same gene serves as a primary indicator of alternative transcript usage.


\section{DTU Fundamentals}

\subsection{DTU vs Differential Gene Expression}
\gls{dtu} occurs when the relative expression of multiple transcripts arising from the same gene changes between different conditions. This differs fundamentally from \gls{deg}, which measures total abundance changes.

Biologically, DTU reflects alternative splicing decisions—a key regulatory mechanism where cells produce functionally distinct protein isoforms from the same gene. Detecting these shifts requires transcript-level resolution; aggregating to genes obliterates the signal.

\subsection{Why Bulk Methods Fail on scRNA-seq}
Established bulk RNA-seq DTU methods\defn{rMATS / SUPPA2 / DRIMSeq}{Bulk DTU tools that model continuous isoform proportions from deep internal read coverage.} (rMATS, SUPPA2, DRIMSeq) model continuous isoform proportions and require deep internal read coverage; they fail on single-cell data for three reasons \cite{galfre2021cotan, love2018swimming}:
\begin{enumerate}
    \item \textbf{Low per-cell depth:} Results in degenerate proportion estimates with excessive zeros.
    \item \textbf{$3'$-bias:} Terminal sequencing renders internal splicing events invisible to junction-spanning methods.
    \item \textbf{Pseudo-bulk aggregation:} Merging cells recovers read depth but obscures the cell-to-cell heterogeneity that drives DTU detection.
\end{enumerate}

Proportion-based DTU methods have been extended to single cells: \textit{satuRn} \cite{gilis2021satuRn} fits a quasi-binomial GLM to isoform proportions, while \textit{Sierra} \cite{patrick2020sierra} and \textit{DTUrtle} \cite{tekath2021dturtle} likewise model isoform proportions, several of them relying on pseudo-bulk aggregation and therefore inheriting the trade-off above. Draper et al.\ \cite{draper2025selecting} survey the practical criteria for selecting differential-splicing methods on short-read data. None of these approaches exploits the zero pattern directly, which is the gap this thesis targets.

This motivates adapting COTAN's sparse-data framework for transcript-level analysis.


\subsection{Mutual Exclusivity Signal}
Within the COTAN framework, negative statistical associations (\gls{mutualExclusivity}) between features can indicate alternative usage patterns. For isoforms of the same gene, such mutual exclusivity would be consistent with alternative splicing or \gls{isoformSwitch} events where different cell states preferentially express distinct transcript variants.

This thesis adapts this principle for DTU detection by screening for mutually exclusive isoform pairs and quantifying their cluster-specific preferences—a methodological contribution detailed in Chapter 3.

